% related-work.tex -- BNR-2026-012

\section{Related Work}
\label{sec:related}

\subsection{Production DAC Architectures}

\textbf{EigenDA V2}~\cite{eigenda2024v2} implements Reed-Solomon erasure
coding with KZG polynomial commitments across 200+ operators. Data is
encoded with $8\times$ redundancy (8{,}192 chunks, any 1{,}024 sufficient
for reconstruction). Write throughput reaches 100\,MB/s with read
throughput of 2{,}000\,MB/s. End-to-end latency averages 5\,s with a
P99 of 10\,s. Blob sizes reach 16\,MB per blob. While EigenDA targets
public rollups with large validator sets, its use of RS+KZG validates
the erasure coding approach at scale.

\textbf{Arbitrum AnyTrust (Nova)}~\cite{anytrust} operates a 6-member
committee with a 5-of-6 threshold for Data Availability Certificates
(DACerts). The protocol uses BLS signature aggregation and provides an
automatic fallback to full rollup mode when the DACert threshold is not
met. The honest minority assumption requires at least 2 of 6 members to
be honest to prevent false attestation. AnyTrust does not employ erasure
coding; each member stores the complete batch data.

\textbf{Polygon CDK DAC}~\cite{polygoncdk} implements a configurable
$M$-of-$N$ ECDSA signature scheme via the \texttt{CDKDataCommittee.sol}
contract. Committee members store full data without erasure coding or
privacy protections. On-chain verification costs 7 $\times$
\texttt{ecrecover} $\approx$ 21{,}000 gas for a 7-member committee.

\textbf{Celestia}~\cite{celestia} employs two-dimensional Reed-Solomon
encoding with a maximum square of $512 \times 512 = 262{,}144$ shares at
512 bytes each ($\approx$128\,MB). Block time is approximately 12
seconds. The LeoRSCodec library provides $O(n \log^2 n)$ encoding via
Leopard Reed-Solomon. While Celestia targets general-purpose data
availability sampling (DAS), its 2D RS approach demonstrates the
scalability of erasure coding at the protocol level.

\textbf{Avail}~\cite{avail} extends the Celestia model on a modified
Substrate chain with KZG commitments per block row. Block time is 20
seconds with a mainnet maximum block size of 2\,MB (scaling to 128\,MB).
KZG commitment generation requires approximately 100\,ms per 512 bytes.

Table~\ref{tab:dac-comparison} summarizes the key architectural
differences.

\begin{table}[t]
\centering
\caption{Production DAC Architecture Comparison}
\label{tab:dac-comparison}
\begin{tabular}{@{}lcccc@{}}
\toprule
\textbf{System} & \textbf{Erasure} & \textbf{Privacy} & \textbf{Nodes} & \textbf{Sigs} \\
\midrule
EigenDA V2  & RS+KZG  & None   & 200+  & BLS    \\
AnyTrust    & None    & None   & 6     & BLS    \\
Polygon CDK & None    & None   & $N$   & ECDSA  \\
Celestia    & 2D RS   & None   & 100+  & EdDSA  \\
Avail       & 2D RS   & None   & 100+  & BABE   \\
RU-V6       & None    & SSS    & 3     & ECDSA  \\
\textbf{This work} & \textbf{RS} & \textbf{AES+SSS} & \textbf{7} & \textbf{ECDSA*} \\
\bottomrule
\multicolumn{5}{l}{\footnotesize *Production target: BLS aggregation.} \\
\end{tabular}
\end{table}

\subsection{Erasure Coding for Data Availability}

Reed-Solomon codes~\cite{reedsolomon1960} provide optimal erasure
correction: a $(k,n)$ code over a finite field encodes $k$ data symbols
into $n$ coded symbols such that any $k$ suffice for reconstruction. The
coding rate $R = k/n$ determines storage efficiency: total storage is
$(n/k)\times$ the original data size.

Al-Bassam et al.~\cite{albassam2018fraud} introduced fraud and data
availability proofs for blockchain light clients, establishing the
theoretical foundation for erasure-coded data availability. Nazirkhanova
et al.~\cite{nazirkhanova2021} proposed Semi-AVID-PR (Asynchronous
Verifiable Information Dispersal with Provable Retrievability), achieving
dispersal of 22\,MB across 256 nodes in under 3 seconds.

Hall-Andersen et al.~\cite{hallandersen2023} formalized the foundations
of data availability sampling, proving security bounds for 2D RS
schemes. Alhaddad et al.~\cite{alhaddad2022} extended verifiable
information dispersal to asynchronous networks, providing the
theoretical basis for DAC protocols that tolerate network partitions.

\subsection{KZG Polynomial Commitments}

Kate, Zaverucha, and Goldberg~\cite{kate2010} introduced constant-size
polynomial commitments enabling verifiable encoding: each RS chunk can be
verified against a single 48-byte commitment without accessing other
chunks. Verification requires two pairing operations on BN254, costing
approximately 124{,}000 gas on EVM via the \texttt{ecPairing}
precompile~\cite{eip196, eip197}. On the zero-fee Basis Network L1,
verification is computationally bounded rather than cost bounded.

\subsection{BLS Signature Aggregation}

Boneh, Drijvers, and Neven~\cite{boneh2018} formalized compact
multi-signatures enabling constant-size attestation regardless of
committee size. For a 7-member committee, BLS aggregation reduces
on-chain signature storage from 455 bytes (7 ECDSA signatures) to 48
bytes (one aggregated G1 point). Verification requires two pairing
operations ($\approx$124{,}000 gas), compared to 21{,}000 gas for 7
\texttt{ecrecover} calls. The trade-off favors BLS at larger committee
sizes where constant-size signatures dominate.

\subsection{Shamir Secret Sharing}

Shamir's $(k,n)$-threshold secret sharing scheme~\cite{shamir1979}
provides information-theoretic security: any coalition of fewer than $k$
shareholders learns zero information about the secret. Our prior work
(RU-V6) applied Shamir sharing directly to batch data, achieving strong
privacy at the cost of $n\times$ storage overhead ($n$ full-size shares)
and $O(k^2)$ Lagrange interpolation per field element during recovery.
This paper restricts Shamir to the 32-byte AES key, reducing the cost to
a single field element operation while preserving information-theoretic
key secrecy.
